% !TeX program = xelatex
% !TeX encoding = UTF-8
% Build: latexmk -xelatex -interaction=nonstopmode -halt-on-error main.tex
% Professional scientific report, version 1.0.0, 15 November 2026.
% Faithfully typeset from paternalizm_trauma_PTSD_CPTSD.md.
\documentclass[11pt,a4paper,oneside,headinclude=true,footinclude=false]{scrartcl}
\usepackage{fontspec}
\usepackage[polish]{babel}
\usepackage{microtype}
\usepackage{geometry}
\usepackage{setspace}
\usepackage[table]{xcolor}
\usepackage{array}
\usepackage{booktabs}
\usepackage{longtable}
\usepackage{calc}
\usepackage{pdflscape}
\usepackage{tikz}
\usetikzlibrary{arrows.meta,positioning,calc}
\usepackage[most]{tcolorbox}
\usepackage{enumitem}
\usepackage{scrlayer-scrpage}
\usepackage{hyperref}
\usepackage{bookmark}
\usepackage{url}
\usepackage{needspace}
\usepackage{ragged2e}

\setmainfont{Noto Serif}[Ligatures=TeX,Numbers=Proportional]
\setsansfont{Inter}[Ligatures=TeX]
\setmonofont{DejaVu Sans Mono}[Scale=0.84]
\geometry{left=27mm,right=25mm,top=24mm,bottom=26mm,headsep=8mm,footskip=10mm}
\setstretch{1.155}
\definecolor{navy}{HTML}{1D344E}
\definecolor{teal}{HTML}{156C79}
\definecolor{ink}{HTML}{273441}
\definecolor{muted}{HTML}{637382}
\definecolor{paleblue}{HTML}{F0F6F8}
\definecolor{tablehead}{HTML}{E8F2F4}
\definecolor{borderline}{HTML}{C7D9DC}
\color{ink}
\hypersetup{
  pdftitle={Paternalizm jako źródło cierpienia psychicznego i czynnik sprzyjający traumatyzacji},
  pdfsubject={Krytyczna analiza psychologiczna, kliniczna i neurobiologiczna},
  pdfkeywords={paternalizm, PTSD, CPTSD, autonomia, trauma, przymus},
  colorlinks=true, linkcolor=navy, urlcolor=teal, citecolor=teal,
  pdfcreator={XeLaTeX}, unicode=true,
  bookmarksopen=true, bookmarksnumbered=true
}
\setkomafont{section}{\sffamily\bfseries\Large\color{navy}}
\setkomafont{subsection}{\sffamily\bfseries\normalsize\color{teal}}
\setkomafont{pageheadfoot}{\sffamily\footnotesize}
\RedeclareSectionCommand[beforeskip=2.0em,afterskip=0.85em]{section}
\RedeclareSectionCommand[beforeskip=1.5em,afterskip=0.6em]{subsection}
\setlength{\parindent}{0pt}
\setlength{\parskip}{5pt plus 1.2pt minus 0.8pt}
\setlist{topsep=5pt,itemsep=3pt,parsep=0pt,leftmargin=1.9em}
\providecommand{\tightlist}{\setlength{\itemsep}{3pt}\setlength{\parskip}{0pt}}
\providecommand{\real}[1]{#1}
\setlength{\LTpre}{0.7em}
\setlength{\LTpost}{1.2em}
\renewcommand{\arraystretch}{1.28}
\AtBeginDocument{\emergencystretch=2em}
\widowpenalty=8000
\clubpenalty=8000
\displaywidowpenalty=8000
\raggedbottom
\urlstyle{same}
\setlength{\emergencystretch}{1.5em}
\newtcolorbox{keyfinding}{
  enhanced,breakable,sharp corners,boxrule=0pt,frame hidden,
  borderline west={2.2pt}{0pt}{teal},colback=paleblue,
  left=11pt,right=11pt,top=9pt,bottom=9pt,before skip=9pt,after skip=10pt
}
\newenvironment{apareferences}{%
  \par\begingroup\small\setlength{\leftskip}{6mm}%
  \setlength{\parindent}{-6mm}%
  \setlength{\parskip}{9pt plus 1pt minus 1pt}%
}{\par\endgroup}
\newtcolorbox{editorialnote}{
  enhanced,breakable,sharp corners,boxrule=0.4pt,
  colframe=borderline,colback=white,
  left=9pt,right=9pt,top=8pt,bottom=8pt,before skip=13pt,after skip=10pt,
  fontupper=\small\color{muted}
}
\clearpairofpagestyles
\ihead{\textcolor{navy}{\textbf{PATERNALIZM}}\textcolor{muted}{\enspace / \enspace PTSD · CPTSD}}
\ohead{\textcolor{muted}{ANALIZA NAUKOWA}}
\ifoot{\textcolor{muted}{Wersja 1.0.0 \enspace · \enspace 15.11.2026}}
\ofoot{\textcolor{muted}{\pagemark}}
\setheadsepline{0.45pt}
\pagestyle{scrheadings}
\KOMAoptions{headsepline=true}

\begin{document}
\begin{titlepage}
\thispagestyle{empty}
\begin{tikzpicture}[remember picture,overlay]
  \fill[navy] (current page.north west) rectangle ([yshift=-10mm]current page.north east);
  \fill[teal] ([xshift=26mm,yshift=-49mm] current page.north west) rectangle ++(14mm,1.3mm);
  \draw[line width=0.7pt,teal!35] ([xshift=-22mm,yshift=-37mm]current page.north east) circle (22mm);
  \draw[line width=0.7pt,teal!20] ([xshift=-22mm,yshift=-37mm]current page.north east) circle (29mm);
  \draw[line width=0.7pt,teal!12] ([xshift=-22mm,yshift=-37mm]current page.north east) circle (36mm);
  \fill[teal] ([xshift=27mm,yshift=32mm]current page.south west) circle (2.1mm);
\end{tikzpicture}
\vspace*{40mm}
{\sffamily\bfseries\fontsize{10}{13}\selectfont\color{teal}\MakeUppercase{Raport naukowy · analiza krytyczna}\par}
\vspace{14mm}
{\sffamily\bfseries\color{navy}\fontsize{26}{32}\selectfont
Paternalizm jako źródło\\[2pt]
cierpienia psychicznego\\[2pt]
i czynnik sprzyjający\\[2pt]
traumatyzacji\par}
\vspace{12mm}
\noindent{\color{teal}\rule{32mm}{1.35pt}}\par
\vspace{8mm}
{\sffamily\fontsize{14}{21}\selectfont\color{ink}
Krytyczna analiza psychologiczna,\\
kliniczna i neurobiologiczna\par}
\vspace{9mm}
{\sffamily\fontsize{10.5}{16}\selectfont\color{muted}
Autonomia · kontrola interpersonalna · PTSD · CPTSD\par}
\vfill
{\sffamily\fontsize{9.4}{13}\selectfont\color{muted}
\textbf{Wersja}\enspace 1.0.0\qquad
\textbf{Data opracowania składu}\enspace 15 listopada 2026\par}
\vspace{2mm}
{\sffamily\fontsize{9.0}{12}\selectfont\color{muted}
\end{titlepage}
\setcounter{page}{1}
\pdfbookmark[1]{Spis treści}{toc}
\tableofcontents
\clearpage

\section{Streszczenie naukowe}\label{sec:1}

\textbf{Paternalizm może wywoływać negatywne konsekwencje psychologiczne
nawet wtedy, gdy jest motywowany autentyczną troską o drugiego
człowieka.} Szczególnie istotne są sytuacje, w których działania
ochronne prowadzą do systematycznego ograniczania autonomii, podważania
kompetencji, narzucania decyzji i utrwalania zależności
interpersonalnej.

Dostępne badania pozwalają wyróżnić trzy zasadnicze drogi, przez które
paternalizm może uczestniczyć w powstawaniu lub podtrzymywaniu
cierpienia psychicznego:

\begin{enumerate}
\def\labelenumi{\arabic{enumi}.}
\tightlist
\item
  \textbf{Droga bezpośredniego obciążenia psychologicznego}
  ograniczanie samostanowienia może prowadzić do frustracji potrzeb
  psychologicznych, poczucia bezradności, wstydu, obniżenia własnej
  skuteczności, lęku i konfliktów interpersonalnych. Nie wszystkie te
  skutki zostały jednak bezpośrednio zbadane w kontekście
  paternalizmu.
\item
  \textbf{Droga kumulatywnej traumatyzacji interpersonalnej}
  paternalizm może stanowić element szerszego wzorca dominacji, kontroli
  przymusowej, izolacji i przemocy psychicznej. W takich okolicznościach
  ograniczanie autonomii przestaje być jedynie stylem sprawowania
  opieki, a może współtworzyć środowisko chronicznego zagrożenia.
\item
  \textbf{Droga podtrzymywania następstw wcześniejszej traumy} u osób,
  które wcześniej doświadczały przymusu, przemocy lub pozbawienia
  kontroli, późniejsze zachowania paternalistyczne mogą przypominać
  wcześniejsze relacje dominacji, wzmacniać negatywne przekonania o
  sobie i utrudniać odzyskiwanie sprawczości.
\end{enumerate}

Najsilniejsze dowody dotyczą nie paternalizmu rozumianego ogólnie, lecz
zjawisk pokrewnych: kontroli przymusowej, przemocy psychicznej,
nadopiekuńczości rodzicielskiej oraz negatywnych reakcji otoczenia na
ujawnienie przemocy.

Przykładowo, metaanaliza Lohmann i współpracowników (2024), obejmująca
45 badań w syntezie ilościowej, wykazała umiarkowany związek kontroli
przymusowej z objawami PTSD (\emph{r} = 0,32; 95\% CI {[}0,28, 0,37{]}).
Jest to istotny dowód na znaczenie chronicznej kontroli
interpersonalnej. Podobnie przegląd Afzal i współpracowników (2023) wskazuje, że
nadopiekuńczość rodziców oraz niektóre sposoby unikania rozmów o traumie
są powiązane z nasileniem objawów pourazowych u dzieci.

\begin{keyfinding}
\textbf{Zasadniczy wniosek:} paternalizm nie musi sam w sobie stanowić
zdarzenia traumatycznego w sensie diagnostycznym. Może jednak być
źródłem rzeczywistego cierpienia psychicznego, współtworzyć
traumatyzujące środowisko interpersonalne albo utrwalać konsekwencje
wcześniejszych doświadczeń traumatycznych.
\end{keyfinding}

W dalszej analizie kluczowe jest rozróżnienie między negatywnym skutkiem
psychologicznym, doświadczeniem subiektywnie traumatycznym oraz
ekspozycją spełniającą kryteria diagnostyczne PTSD lub CPTSD.

\section{Paternalizm, a doświadczenie traumatyczne}\label{sec:2}

Paternalizm oznacza zastępowanie lub ograniczanie decyzji innej osoby w
przekonaniu, że służy to jej dobru. Jego istotą nie musi być wrogość.
Może wynikać z troski, lęku o bezpieczeństwo, przekonania o własnych
kompetencjach lub poczucia odpowiedzialności.

Psychologiczna szkodliwość paternalizmu zależy jednak nie tylko od
intencji osoby sprawującej kontrolę, ale również od zakresu ograniczeń,
możliwości sprzeciwu, asymetrii władzy i konsekwencji dla odbiorcy.

Należy odróżnić:

\begin{itemize}
\tightlist
\item
  \textbf{Paternalizm życzliwy:} ograniczanie samodzielności w celu
  ochrony, bez zamiaru wyrządzenia krzywdy.
\item
  \textbf{Paternalizm kontrolujący:} narzucanie decyzji i sposobów
  postępowania mimo zachowanej zdolności osoby do samostanowienia.
\item
  \textbf{Paternalizm przymusowy:} egzekwowanie decyzji przez sankcje,
  ograniczenia swobody lub inne formy przymusu.
\item
  \textbf{Kontrolę przymusową (\emph{coercive control}):} utrwalony
  wzorzec dominacji, ograniczania wolności i podporządkowania. Może
  wykorzystywać retorykę troski, ale nie jest synonimem paternalizmu.
\end{itemize}

Równie ważne jest rozróżnienie pojęć klinicznych. W DSM-5-TR rozpoznanie
PTSD wymaga określonego rodzaju ekspozycji, obejmującego rzeczywistą lub
grożącą śmierć, poważne obrażenia albo przemoc seksualną. ICD-11 odnosi
PTSD i CPTSD do ekspozycji na zdarzenie lub serię zdarzeń skrajnie
zagrażających bądź przerażających.

Przewlekłe upokarzanie, infantylizacja i ograniczanie autonomii mogą
powodować istotne cierpienie psychiczne, nawet jeżeli nie spełniają
kryteriów ekspozycji traumatycznej wymaganych do rozpoznania PTSD.

\section{Negatywne konsekwencje psychologiczne paternalizmu}\label{sec:3}

\subsection{Ograniczenie autonomii i frustracja podstawowych potrzeb psychologicznych}\label{subsec:3.1}

Według teorii samostanowienia autonomia, kompetencja i relacyjność
należą do podstawowych potrzeb psychologicznych. Paternalizm może
oddziaływać negatywnie na każdą z nich.

Osoba poddawana nadmiernej kontroli może doświadczać:

\begin{itemize}
\tightlist
\item
  ograniczenia poczucia wpływu na własne życie;
\item
  obniżenia motywacji do samodzielnego działania;
\item
  frustracji, złości i reaktancji;
\item
  poczucia niekompetencji;
\item
  narastającej zależności od innych;
\item
  osłabienia zaufania do własnych ocen.
\end{itemize}

W trzech badaniach Bartholomew et al.~(2011) wykazali, że
interpersonalny styl kontrolujący wiązał się z frustracją potrzeb
psychologicznych, a ta z negatywnymi wskaźnikami funkcjonowania, w tym
depresyjnością, wypaleniem i negatywnym afektem.

\textbf{Znaczenie dla traumy:} jeżeli jednostka wcześniej doświadczyła
przemocy związanej z utratą kontroli, ponowne ograniczenie autonomii
może nie tylko wywołać bieżący stres, lecz również potwierdzać
przekonanie, że nie ma ona wpływu na własne bezpieczeństwo.

\subsection{Podważanie kompetencji i poczucia własnej skuteczności}\label{subsec:3.2}

Paternalizm może komunikować odbiorcy, że jest mniej kompetentny, niż
rzeczywiście jest.

Przykładowo, powtarzane stwierdzenia:

\begin{quote}
„Nie podejmuj tej decyzji. Ja wiem lepiej, ponieważ ty nie potrafisz
właściwie ocenić sytuacji''.
\end{quote}

mogą stopniowo osłabiać zaufanie jednostki do własnych zdolności
poznawczych i decyzyjnych.

Istnieją bezpośrednie dowody eksperymentalne na negatywne efekty
określonych form paternalizmu życzliwego.

Dardenne et al.~(2007) w czterech eksperymentach wykazali, że ekspozycja
na życzliwy seksizm, obejmujący paternalistyczne przekazy o
kompetencjach kobiet, pogarszała ich wykonanie zadań poznawczych. W
jednym z eksperymentów efekt był mediowany przez intruzywne myśli
dotyczące własnej kompetencji. Autorzy odróżnili przy tym skutki
życzliwego seksizmu od samego otrzymywania nieproszonej pomocy.

Wyniki te pokazują, że komunikat o pozornie pozytywnym charakterze może
wywoływać mierzalne negatywne efekty poznawcze.

Nie dowodzą jednak, że takie oddziaływanie samo w sobie powoduje traumę.

\subsection{Wyuczona bezradność i utrwalanie zależności}\label{subsec:3.3}

Powtarzające się przejmowanie decyzji za drugą osobę może ograniczać jej
okazje do samodzielnego działania, ponoszenia odpowiedzialności i
doświadczania skuteczności własnych wyborów.

Potencjalny mechanizm obejmuje następującą sekwencję:

\begin{enumerate}
\def\labelenumi{\arabic{enumi}.}
\tightlist
\item
  Powtarzające się przejmowanie decyzji.
\item
  Ograniczenie doświadczeń samodzielnej skuteczności.
\item
  Osłabienie przekonania o możliwości wpływania na sytuację.
\item
  Bierność, zależność lub unikanie decyzji.
\end{enumerate}

Jest to model hipotetyczny, a nie empirycznie potwierdzona sekwencja
przyczynowa specyficzna dla paternalizmu.

Nygaard et al.~(2017) wykazali w badaniu podłużnym osób po napaści
fizycznej, że poczucie własnej skuteczności i objawy pourazowe pozostają
ze sobą w relacji zmieniającej się w czasie. Wyniki sugerowały, że
objawy pourazowe mogą początkowo obniżać własną skuteczność, natomiast
większa skuteczność może sprzyjać późniejszemu zdrowieniu.

To ważny argument za badaniem autonomii jako potencjalnego mechanizmu
zdrowienia, lecz nie dowód, że paternalizm wywołuje wyuczoną bezradność.

\subsection{Infantylizacja, wstyd i negatywny obraz siebie}\label{subsec:3.4}

Infantylizacja polega na traktowaniu osoby jako mniej dojrzałej lub
mniej kompetentnej, niż wynikałoby to z jej rzeczywistych możliwości.

Może przyjmować formę nadmiernego pouczania, odbierania
odpowiedzialności, pomijania opinii albo komunikowania, że jednostka nie
jest zdolna do podejmowania decyzji.

Potencjalnymi konsekwencjami są:

\begin{itemize}
\tightlist
\item
  obniżenie samooceny;
\item
  wstyd związany z zależnością;
\item
  wycofanie z aktywności;
\item
  niechęć do podejmowania inicjatywy;
\item
  osłabienie poczucia dorosłości i podmiotowości.
\end{itemize}

Badanie Hammond i Overall (2015), obejmujące obserwację interakcji 100
par, wykazało, że mężczyźni silniej akceptujący życzliwy seksizm
częściej udzielali partnerkom pomocy zorientowanej na zależność,
polegającej m.in. na dostarczaniu gotowych rozwiązań z pominięciem ich
kompetencji. Wiązało się to z niższym poczuciem kompetencji i mniejszym
poczuciem bycia pozytywnie ocenianą przez partnera.

Badanie dostarcza bezpośrednich dowodów na psychologiczne koszty
paternalistycznego wsparcia w relacji partnerskiej. Nie oceniano jednak
rozwoju PTSD.

\section{Jak paternalizm może prowadzić do doświadczeń traumatycznych w ciągu życia?}\label{sec:4}

\subsection{Dzieciństwo: nadopiekuńczość, kontrola i ograniczenie rozwoju sprawczości}\label{subsec:4.1}

Dzieciństwo jest okresem szczególnie istotnym dla kształtowania poczucia
kontroli, autonomii oraz przekonań dotyczących własnych kompetencji.

Rodzicielski paternalizm może stać się problematyczny, gdy ochrona jest
nieproporcjonalna do wieku, możliwości dziecka i rzeczywistego ryzyka.

Przykłady obejmują systematyczne uniemożliwianie samodzielnego
rozwiązywania problemów, nadmierne monitorowanie, ograniczanie kontaktów
społecznych oraz podejmowanie wszystkich decyzji za dziecko.

Przegląd Vigdal i Brønnick (2022) wykazał powiązania tzw.
\emph{helicopter parenting} z objawami lękowymi i depresyjnymi. Autorzy
zaznaczyli jednak, że jakość dostępnych badań podłużnych nie pozwalała
jednoznacznie ustalić kierunku przyczynowego.

Szczególnie ważne są badania dzieci, które doświadczyły już zdarzeń
traumatycznych.

Bokszczanin (2008), analizując 533 uczniów 28 miesięcy po powodzi w
Polsce, wykazała, że nadopiekuńczość rodzicielska przewidywała nasilenie
objawów PTSD i moderowała związek między ekspozycją na katastrofę a
objawami pourazowymi.

Hiller et al.~(2018) dostarczyli ponadto dowodów podłużnych
wskazujących, że niektóre reakcje rodziców po traumie mogą oddziaływać
na objawy dzieci poprzez ich interpretacje poznawcze i strategie
radzenia sobie. Nie stwierdzono natomiast jednoznacznego ochronnego
wpływu wszystkich form pozytywnego wsparcia rodzicielskiego.

\textbf{Interpretacja:} nadmierna ochrona może ograniczać proces
odzyskiwania poczucia bezpieczeństwa poprzez własne działanie. Nie
oznacza to jednak, że sama nadopiekuńczość jest wystarczającą przyczyną
PTSD.

\subsection{Relacje partnerskie: przejście od troski do kontroli przymusowej}\label{subsec:4.2}

Jednym z najistotniejszych zagrożeń jest możliwość wykorzystania
paternalistycznego uzasadnienia do stopniowego rozszerzania kontroli nad
partnerem.

Przykładowa sekwencja może wyglądać następująco:

\begin{enumerate}
\def\labelenumi{\arabic{enumi}.}
\tightlist
\item
  Partner deklaruje troskę i potrzebę ochrony.
\item
  Zaczyna podejmować decyzje dotyczące kontaktów, finansów lub
  aktywności drugiej osoby.
\item
  Ograniczenia stają się coraz bardziej rozległe.
\item
  Próby odzyskania autonomii spotykają się z naciskiem, wzbudzaniem winy
  lub sankcjami.
\item
  Osoba kontrolowana stopniowo traci możliwość swobodnego opuszczenia
  relacji.
\end{enumerate}

Nie każda paternalistyczna relacja rozwija się w ten sposób. Jest to
jednak możliwy mechanizm przejścia od ochrony do dominacji.

W badaniu podłużnym Taft et al.~(2005) przemoc psychiczna była silniej i
bardziej niezależnie powiązana z objawami PTSD niż przemoc fizyczna.
Szczególnie istotne były poniżanie, restrykcyjne ograniczanie autonomii
partnerki oraz dominacja i zastraszanie.

Metaanaliza Lohmann et al.~(2024) potwierdza umiarkowany związek
kontroli przymusowej z PTSD. Jednocześnie autorzy wskazali, że dane
dotyczące CPTSD były niewystarczające do przeprowadzenia odrębnej
metaanalizy --- tylko jedno badanie dotyczyło bezpośrednio tego związku.

W tym kontekście paternalizm może pełnić dwie różne funkcje:
rzeczywistego motywu działania opiekuńczego albo uzasadnienia dla
praktyk kontrolujących. Nie można ich utożsamiać bez analizy zachowania
i jego konsekwencji.

\subsection{Paternalizm epistemiczny i wtórna wiktymizacja}\label{subsec:4.3}

Paternalizm epistemiczny występuje wtedy, gdy perspektywa danej osoby
jest systematycznie pomijana lub zastępowana interpretacją osoby
posiadającej większy autorytet.

W kontekście traumy może przyjmować postać stwierdzeń:

\begin{quote}
„Źle pamiętasz to wydarzenie''.
\end{quote}

\begin{quote}
„Twoja trauma sprawia, że nie potrafisz ocenić sytuacji''.
\end{quote}

\begin{quote}
„Wiem lepiej od ciebie, co naprawdę przeżywasz''.
\end{quote}

Takie komunikaty nie muszą być nieprawdziwe w każdej sytuacji. Problem
pojawia się, gdy są formułowane bez wystarczających podstaw, z
pominięciem wyjaśnień i bez uznania możliwości przedstawienia odmiennej
perspektywy.

W badaniu Ullman i Peter-Hagene (2014), obejmującym 1863 osoby po
przemocy seksualnej, negatywne reakcje otoczenia na ujawnienie przemocy
wiązały się z większym nasileniem PTSD, zarówno bezpośrednio w modelu
statystycznym, jak i pośrednio przez nieadaptacyjne strategie radzenia
sobie. Pozytywne reakcje wiązały się z większym poczuciem kontroli nad
zdrowieniem, które z kolei było powiązane z mniejszym nasileniem
objawów.

Badanie nie mierzyło paternalizmu epistemicznego jako odrębnego
konstruktu. Pokazuje jednak, że sposób reagowania na doświadczenie osoby
pokrzywdzonej ma znaczenie dla jej funkcjonowania pourazowego.

\subsection{Instytucje: zdrada zaufania i ponowne doświadczenie bezsilności}\label{subsec:4.4}

Paternalizm instytucjonalny może być szczególnie obciążający, ponieważ
jednostka często ma ograniczoną możliwość zmiany instytucji,
zakwestionowania decyzji lub uniknięcia konsekwencji odmowy.

Dotyczy to między innymi szpitali, placówek opiekuńczych, instytucji
edukacyjnych i systemów pomocy społecznej.

Pokrewnym konstruktem jest \emph{institutional betrayal}, czyli zdrada
instytucjonalna: sytuacja, w której instytucja odpowiedzialna za
bezpieczeństwo zawodzi osobę pokrzywdzoną.

Smith i Freyd (2013) wykazały, że zdrada instytucjonalna towarzysząca
doświadczeniom przemocy seksualnej wiązała się z większym nasileniem
lęku, dysocjacji i innych negatywnych następstw psychologicznych.

Paternalizm i zdrada instytucjonalna nie są tożsame. Mogą jednak
współwystępować, np. gdy instytucja ogranicza możliwość zgłoszenia
przemocy, uzasadniając to ochroną osoby pokrzywdzonej.

\subsection{Opieka psychiatryczna: paternalizm i przymus}\label{subsec:4.5}

Najbardziej bezpośrednie dowody na traumatyzujący potencjał działań
ochronnych pochodzą z badań nad przymusem psychiatrycznym.

Guzmán-Parra et al.~(2019) badali 111 pacjentów hospitalizowanych
psychiatrycznie, wobec których zastosowano przymusową farmakoterapię,
unieruchomienie mechaniczne albo obie interwencje. Unieruchomienie
mechaniczne oraz połączenie obu środków wiązały się z wyższymi
wskaźnikami stresu pourazowego niż sama przymusowa farmakoterapia.

Badanie nie pozwala ustalić, czy różnice wynikały bezpośrednio z
interwencji, wcześniejszego stanu pacjentów czy innych czynników.

Pokazuje jednak, że \textbf{działanie uzasadniane bezpieczeństwem
pacjenta może samo być źródłem intensywnego stresu i negatywnego
doświadczenia psychicznego}.

W tym przypadku źródłem zagrożenia jest przede wszystkim zastosowany
przymus, a nie sama deklarowana troska.

\section{Paternalizm a PTSD i CPTSD}\label{sec:5}

W PTSD istotne są między innymi ponowne przeżywanie traumy, unikanie i
utrzymujące się poczucie aktualnego zagrożenia. CPTSD według ICD-11
obejmuje dodatkowo trzy domeny zaburzeń organizacji Ja (DSO).

Paternalizm może potencjalnie oddziaływać na każdą z nich, ale stopień
potwierdzenia empirycznego jest różny.

\begingroup\small\setlength{\tabcolsep}{5pt}
\begin{longtable}[]{@{}
  >{\raggedright\arraybackslash}p{(\columnwidth - 4\tabcolsep) * \real{0.3333}}
  >{\raggedright\arraybackslash}p{(\columnwidth - 4\tabcolsep) * \real{0.3333}}
  >{\raggedright\arraybackslash}p{(\columnwidth - 4\tabcolsep) * \real{0.3333}}@{}}
\toprule\noalign{}
\rowcolor{tablehead}
\begin{minipage}[b]{\linewidth}\raggedright
Domena
\end{minipage} & \begin{minipage}[b]{\linewidth}\raggedright
Możliwy wpływ paternalizmu
\end{minipage} & \begin{minipage}[b]{\linewidth}\raggedright
Status dowodów
\end{minipage} \\
\midrule\noalign{}
\endhead
\bottomrule\noalign{}
\endlastfoot
Dysregulacja afektu & Lęk, złość, wstyd i trudności regulacyjne po
utracie kontroli & Dowody pośrednie \\
Negatywny obraz siebie & Utrwalanie przekonań o niekompetencji,
bezradności i zależności & Dowody pośrednie i eksperymentalne poza
CPTSD \\
Zaburzenia relacji & Nieufność, podporządkowanie, unikanie bliskości lub
nadmierna zależność & Głównie dowody pośrednie i hipotezy \\
\end{longtable}
\endgroup

W szczególności należy rozważyć, czy paternalizm może utrwalać
charakterystyczne dla niektórych osób z CPTSD przekonania: „nie jestem
zdolny do samodzielnego działania'', „inni mają nade mną władzę'' albo
„nie mogę ufać własnym decyzjom''.

Nie ma jednak wystarczających danych, aby stwierdzić, że osoby z CPTSD
są zawsze bardziej wrażliwe na paternalizm niż osoby z PTSD lub bez
rozpoznania zaburzeń pourazowych.

Istotne może być raczej podobieństwo bieżącej sytuacji do wcześniejszego
doświadczenia przemocy, indywidualna interpretacja ograniczeń i aktualna
możliwość odmowy.

\section{Mechanizmy neuropsychologiczne}\label{sec:6}

Hipoteza neuropsychologiczna zakłada, że zachowania paternalistyczne
mogą u niektórych osób po traumie działać jako bodźce przypominające
wcześniejszą utratę kontroli.

Można wyróżnić cztery potencjalne mechanizmy.

\textbf{Uczenie asocjacyjne.} Jeżeli w przeszłości kontrola ze strony
opiekuna współwystępowała z zagrożeniem, późniejsze sygnały dominacji
mogą aktywować wyuczone skojarzenia. Nie wymaga to świadomego uznania
bieżącej sytuacji za niebezpieczną.

\textbf{Ocena aktualnego zagrożenia.} Paternalistyczna interakcja może
wzmacniać przekonania o braku kontroli i niemożności skutecznego
działania. Jest to zgodne z poznawczymi modelami podtrzymywania PTSD.

\textbf{Regulacja emocji.} Utrata wyboru może zwiększać pobudzenie
emocjonalne, a u niektórych osób prowadzić do wycofania lub zachowań
uległościowych.

\textbf{Pamięć autobiograficzna.} Podobieństwo interpersonalne między
bieżącą sytuacją a wcześniejszą przemocą może sprzyjać przywoływaniu
niechcianych wspomnień.

Nie należy jednak twierdzić, że paternalizm bezpośrednio aktywuje ciało
migdałowate, zaburza funkcję kory przedczołowej czy powoduje
dysregulację osi HPA u osób z PTSD. Takie szczegółowe mechanizmy
wymagają bezpośrednich badań eksperymentalnych.

\section{Integracyjny model negatywnego wpływu paternalizmu}\label{sec:7}

\textbf{Model hipotetyczny:}

\begin{enumerate}
\def\labelenumi{\arabic{enumi}.}
\tightlist
\item
  \textbf{Paternalistyczne zachowanie:} narzucanie decyzji,
  nadopiekuńczość, podważanie kompetencji, kontrola.
\item
  \textbf{Dwa możliwe mechanizmy pośredniczące:}

  \begin{itemize}
  \tightlist
  \item
    ograniczenie autonomii \(\rightarrow\) spadek sprawczości i poczucia
    kompetencji;\\
  \item
    ocena zagrożenia \(\rightarrow\) skojarzenia z wcześniejszym przymusem.\\
  \end{itemize}
\item
  \textbf{Reakcje psychologiczne:} lęk, wstyd, bezradność, reaktancja,
  unikanie, wycofanie.
\item
  \textbf{Potencjalne konsekwencje:}

  \begin{itemize}
  \tightlist
  \item
    ogólne cierpienie psychiczne: stres, obniżony dobrostan, trudności
    relacyjne;\\
  \item
    możliwe podtrzymywanie objawów pourazowych: PTSD, DSO w CPTSD,
    dysocjacja.
  \end{itemize}
\end{enumerate}


\begin{figure}[htbp]
\centering
\begin{tikzpicture}[>=Latex, node distance=11mm,
  modelnode/.style={draw=borderline, rounded corners=2mm, fill=paleblue,
    align=center, inner xsep=8pt,inner ysep=8pt, font=\footnotesize\sffamily,
    text=ink, minimum height=11mm},
  modelarrow/.style={->,very thick,draw=teal}]
  \node[modelnode,text width=128mm,fill=navy,text=white,draw=navy] (s) at (0,0)
    {\textbf{Zachowania paternalistyczne}\\narzucanie decyzji, nadopiekuńczość, podważanie kompetencji, kontrola};
  \node[modelnode,text width=56mm] (a) at (-35mm,-23mm)
    {\textbf{Ograniczenie autonomii}\\spadek sprawczości i poczucia kompetencji};
  \node[modelnode,text width=56mm] (z) at (35mm,-23mm)
    {\textbf{Ocena zagrożenia}\\skojarzenia z wcześniejszym przymusem};
  \node[modelnode,text width=128mm] (r) at (0,-47mm)
    {\textbf{Reakcje psychologiczne}\\lęk, wstyd, bezradność, reaktancja, unikanie, wycofanie};
  \node[modelnode,text width=128mm] (c) at (0,-69mm)
    {\textbf{Potencjalne konsekwencje}\\ogólne cierpienie psychiczne; możliwe podtrzymywanie objawów pourazowych};
  \draw[modelarrow] (s.south) -- (a.north);
  \draw[modelarrow] (s.south) -- (z.north);
  \draw[modelarrow] (a.south) -- (r.north);
  \draw[modelarrow] (z.south) -- (r.north);
  \draw[modelarrow] (r.south) -- (c.north);
\end{tikzpicture}
\caption{Schemat modelu hipotetycznego. Strzałki ilustrują przewidywane zależności, a nie potwierdzoną przyczynowość.}
\label{fig:model}
\end{figure}

\emph{Autorski model hipotetyczny. Strzałki i zależności oznaczają
przewidywane mechanizmy, nie potwierdzone w całości relacje przyczynowe.
Wpływ mogą modyfikować historia traumy, możliwość odmowy, rzeczywiste
zagrożenie, jakość relacji i kontekst kulturowy.}

Model ten pozwala również rozróżnić dwa rodzaje skutków. Pierwszy to
cierpienie wywoływane przez paternalizm niezależnie od wcześniejszej
traumy. Drugi to możliwe wzmacnianie objawów pourazowych, kiedy
paternalizm występuje u osoby z wcześniejszymi doświadczeniami przemocy
i utraty kontroli.

\clearpage\begin{landscape}
\section{Zestawienie najważniejszych dowodów}\label{sec:8}

\begingroup\small\setlength{\tabcolsep}{5pt}
\begin{longtable}[]{@{}
  >{\raggedright\arraybackslash}p{(\linewidth - 6\tabcolsep) * \real{0.19}}
  >{\raggedright\arraybackslash}p{(\linewidth - 6\tabcolsep) * \real{0.23}}
  >{\raggedright\arraybackslash}p{(\linewidth - 6\tabcolsep) * \real{0.30}}
  >{\raggedright\arraybackslash}p{(\linewidth - 6\tabcolsep) * \real{0.28}}@{}}
\toprule\noalign{}
\rowcolor{tablehead}
\begin{minipage}[b]{\linewidth}\raggedright
Zjawisko
\end{minipage} & \begin{minipage}[b]{\linewidth}\raggedright
Populacja i projekt
\end{minipage} & \begin{minipage}[b]{\linewidth}\raggedright
Główne ustalenie
\end{minipage} & \begin{minipage}[b]{\linewidth}\raggedright
Ocena
\end{minipage} \\
\midrule\noalign{}
\endhead
\bottomrule\noalign{}
\endlastfoot
Życzliwy seksizm i paternalizm & Dardenne et al.~(2007); cztery
eksperymenty & Pogorszenie wykonania poznawczego; rola intruzji
dotyczących kompetencji & Bezpośrednie dla określonej formy
paternalizmu, nie dla traumy \\
Paternalistyczne wsparcie partnerskie & Hammond \& Overall (2015);
obserwacja 100 par & Wsparcie zorientowane na zależność powiązane z
niższym poczuciem kompetencji & Bezpośrednie, ograniczone do badanego
kontekstu \\
Nadopiekuńczość po traumie & Afzal et al.~(2023); przegląd 27 prac &
Związek niektórych reakcji rodzicielskich z objawami PTSS & Umiarkowane,
głównie obserwacyjne \\
Nadopiekuńczość po katastrofie & Bokszczanin (2008); 533 uczniów &
Nasilenie związku między ekspozycją na katastrofę a PTSD &
Obserwacyjne \\
Kontrola przymusowa & Lohmann et al.~(2024); metaanaliza & \emph{r} =
0,32 dla związku z PTSD & Umiarkowane dla korelacji; przyczynowość
nieustalona \\
Przemoc psychiczna & Taft et al.~(2005); badanie podłużne & Dominacja,
poniżanie i restrykcyjna kontrola powiązane z PTSD & Pośrednie dla
paternalizmu \\
Negatywne reakcje na ujawnienie przemocy & Ullman \& Peter-Hagene
(2014); 1863 osoby & Większe nasilenie PTSD, częściowo poprzez strategie
radzenia sobie & Pośrednie; modele mediacyjne nie dowodzą
przyczynowości \\
Zdrada instytucjonalna & Smith \& Freyd (2013); badanie obserwacyjne &
Większy lęk, dysocjacja i inne następstwa traumy & Pośrednie \\
Przymus psychiatryczny & Guzmán-Parra et al.~(2019); 111 pacjentów &
Większy stres pourazowy przy unieruchomieniu niż przy samej przymusowej
farmakoterapii & Bezpośrednie dla przymusu, nie dla paternalizmu jako
takiego \\
\end{longtable}
\endgroup
\end{landscape}\clearpage


\section{Kontrargumenty i granice interpretacji}\label{sec:9}

Koncentracja na szkodach paternalizmu nie oznacza, że każda interwencja
ograniczająca autonomię jest nieuzasadniona.

W sytuacji bezpośredniego zagrożenia życia, poważnego ograniczenia
zdolności decyzyjnej lub konieczności zapobieżenia ciężkiej szkodzie
czasowa interwencja ochronna może być klinicznie i etycznie uzasadniona.
Powinna jednak być proporcjonalna, możliwie najmniej restrykcyjna i
podlegać kontroli.

Istnieją również ważne alternatywne wyjaśnienia obserwowanych
zależności. Osoby z większym nasileniem objawów mogą częściej otrzymywać
pomoc kontrolującą. Nadopiekuńczość może współwystępować z lękiem
opiekuna, konfliktem rodzinnym albo ogólnie niską jakością relacji.
Część odbiorców może też preferować czasowe delegowanie decyzji,
zwłaszcza jeśli zachowuje możliwość wycofania zgody.

Nie można zatem automatycznie interpretować każdej niechcianej pomocy
jako przemocy, ani każdej reakcji stresowej jako retraumatyzacji.

\section{Hipotezy do dalszej weryfikacji}\label{sec:10}

\begingroup\small\setlength{\tabcolsep}{5pt}
\begin{longtable}[]{@{}
  >{\raggedright\arraybackslash}p{(\columnwidth - 4\tabcolsep) * \real{0.3333}}
  >{\raggedright\arraybackslash}p{(\columnwidth - 4\tabcolsep) * \real{0.3333}}
  >{\raggedright\arraybackslash}p{(\columnwidth - 4\tabcolsep) * \real{0.3333}}@{}}
\toprule\noalign{}
\rowcolor{tablehead}
\begin{minipage}[b]{\linewidth}\raggedright
Hipoteza
\end{minipage} & \begin{minipage}[b]{\linewidth}\raggedright
Przewidywanie
\end{minipage} & \begin{minipage}[b]{\linewidth}\raggedright
Kryterium podważające hipotezę
\end{minipage} \\
\midrule\noalign{}
\endhead
\bottomrule\noalign{}
\endlastfoot
H1. Paternalizm obniża poczucie kompetencji & Komunikaty
paternalistyczne powodują większy spadek samooceny kompetencji niż
komunikaty wspierające autonomię & Brak różnicy w eksperymencie o
odpowiedniej mocy, przy przedziale ufności wykluczającym efekt o
znaczeniu praktycznym \\
H2. Historia kontroli przymusowej moderuje reakcję & Osoby z taką
historią silniej reagują na paternalistyczne komunikaty & Brak
przewidywanej interakcji w badaniu o odpowiedniej mocy \\
H3. Autonomia mediuje wpływ na objawy & Zmniejszenie autonomii poprzedza
wzrost objawów pourazowych & Brak przewidywanej sekwencji czasowej lub
efektu pośredniego \\
H4. Paternalizm nasila DSO w CPTSD & Silniejsze związki z negatywnym
obrazem siebie i relacjami niż z podstawowymi objawami PTSD & Równoważne
lub odwrotne zależności po kontroli zmiennych zakłócających \\
H5. Możliwość odmowy ogranicza szkodliwość & Ten sam komunikat wywołuje
mniejszy stres, gdy osoba może rzeczywiście odmówić & Brak efektu
możliwości odmowy przy potwierdzonej skuteczności manipulacji
eksperymentalnej \\
\end{longtable}
\endgroup

Do badań należałoby wykorzystać m.in. \emph{International Trauma
Questionnaire} (ITQ), PCL-5, miary frustracji potrzeb psychologicznych i
narzędzia oceniające postrzeganą kontrolę interpersonalną. Konieczne
jest odróżnienie rzeczywistego ograniczenia autonomii od jego
subiektywnej interpretacji.

\section{Wnioski końcowe}\label{sec:11}

\textbf{Paternalizm może być szkodliwy psychologicznie bez złych
intencji osoby, która go stosuje.} Badania pokazują, że niektóre formy
paternalistycznego wsparcia obniżają poczucie kompetencji, a
nadopiekuńczość i kontrola interpersonalna mogą współwystępować z
gorszym funkcjonowaniem psychicznym.

Związek z traumą należy jednak rozpatrywać na kilku poziomach.

\textbf{Bezpośrednio potwierdzone empirycznie:} określone formy
życzliwego paternalizmu mogą pogarszać funkcjonowanie poznawcze i
poczucie kompetencji. Przymusowe interwencje ochronne mogą wiązać się z
intensywnym stresem pourazowym. Nie oznacza to jednak, że paternalizm
samodzielnie wywołuje PTSD.

\textbf{Potwierdzone pośrednio:} kontrola przymusowa, przemoc
psychiczna, nadopiekuńczość po traumie i negatywne reakcje na ujawnienie
przemocy są związane z nasileniem objawów pourazowych. Nie wolno
przenosić wielkości tych efektów bezpośrednio na życzliwy paternalizm.

\textbf{Hipoteza zgodna z istniejącymi modelami:} u osób z
wcześniejszymi doświadczeniami przemocy paternalizm może odtwarzać
psychologiczne wzorce utraty kontroli, bezradności i podporządkowania,
utrudniając odzyskiwanie sprawczości.

\textbf{Niewystarczająco zbadane:} czy paternalizm jest niezależnym
czynnikiem ryzyka PTSD lub CPTSD, jakie mechanizmy neurobiologiczne
uczestniczą w tej zależności oraz czy osoby z CPTSD są bardziej wrażliwe
na paternalizm niż inne grupy.

Najważniejsze rozróżnienie dotyczy paternalizmu jako samodzielnego
źródła stresu oraz paternalizmu jako elementu szerszego środowiska
przemocy, dominacji lub przymusu. To drugie zjawisko ma znacznie
silniejsze powiązania empiryczne z objawami pourazowymi.


\section{Bibliografia --- APA 7}\label{sec:12}

Poniższe pozycje mają identyfikowalne DOI. Zastosowano cytowania
autor--rok, pełne listy autorów, tytuły artykułów zapisane zgodnie z APA
7 oraz identyfikatory DOI w formie adresów URL.

\begin{apareferences}
Afzal, N., Ye, S., Page, A. C., Trickey, D., Lyttle, M. D., Hiller, R.
M., \& Halligan, S. L. (2023). A systematic literature review of the
relationship between parenting responses and child post-traumatic stress
symptoms. \emph{European Journal of Psychotraumatology, 14}(1), Article
2156053. \url{https://doi.org/10.1080/20008066.2022.2156053}

Bartholomew, K. J., Ntoumanis, N., Ryan, R. M., Bosch, J. A., \&
Thøgersen-Ntoumani, C. (2011). Self-determination theory and diminished
functioning: The role of interpersonal control and psychological need
thwarting. \emph{Personality and Social Psychology Bulletin, 37}(11),
1459--1473. \url{https://doi.org/10.1177/0146167211413125}

Bokszczanin, A. (2008). Parental support, family conflict, and
overprotectiveness: Predicting PTSD symptom levels of adolescents 28
months after a natural disaster. \emph{Anxiety, Stress \& Coping,
21}(4), 325--335. \url{https://doi.org/10.1080/10615800801950584}

Dardenne, B., Dumont, M., \& Bollier, T. (2007). Insidious dangers of
benevolent sexism: Consequences for women's performance. \emph{Journal
of Personality and Social Psychology, 93}(5), 764--779.
\url{https://doi.org/10.1037/0022-3514.93.5.764}

Guzmán-Parra, J., Aguilera-Serrano, C., García-Sanchez, J. A.,
García-Spínola, E., Torres-Campos, D., Villagrán, J. M., Moreno-Küstner,
B., \& Mayoral-Cleries, F. (2019). Experience coercion, post-traumatic
stress, and satisfaction with treatment associated with different
coercive measures during psychiatric hospitalization.
\emph{International Journal of Mental Health Nursing, 28}(2), 448--456.
\url{https://doi.org/10.1111/inm.12546}

Hammond, M. D., \& Overall, N. C. (2015). Benevolent sexism and support
of romantic partner's goals: Undermining women's competence while
fulfilling men's intimacy needs. \emph{Personality and Social Psychology
Bulletin, 41}(9), 1180--1194.
\url{https://doi.org/10.1177/0146167215593492}

Hiller, R. M., Meiser-Stedman, R., Lobo, S., Creswell, C., Fearon, P.,
Ehlers, A., Murray, L., \& Halligan, S. L. (2018). A longitudinal
investigation of the role of parental responses in predicting children's
post-traumatic distress. \emph{Journal of Child Psychology and
Psychiatry, 59}(7), 781--789. \url{https://doi.org/10.1111/jcpp.12846}

Lohmann, S., Cowlishaw, S., Ney, L. J., O'Donnell, M., \& Felmingham, K.
L. (2024). The trauma and mental health impacts of coercive control: A
systematic review and meta-analysis. \emph{Trauma, Violence, \& Abuse,
25}(1), 630--647. \url{https://doi.org/10.1177/15248380231162972}

Nygaard, E., Johansen, V. A., Siqveland, J., Hussain, A., \& Heir, T.
(2017). Longitudinal relationship between self-efficacy and
posttraumatic stress symptoms 8 years after a violent assault: An
autoregressive cross-lagged model. \emph{Frontiers in Psychology, 8},
Article 913. \url{https://doi.org/10.3389/fpsyg.2017.00913}

Smith, C. P., \& Freyd, J. J. (2013). Dangerous safe havens:
Institutional betrayal exacerbates sexual trauma. \emph{Journal of
Traumatic Stress, 26}(1), 119--124.
\url{https://doi.org/10.1002/jts.21778}

Taft, C. T., Murphy, C. M., King, L. A., Dedeyn, J. M., \& Musser, P. H.
(2005). Posttraumatic stress disorder symptomatology among partners of
men in treatment for relationship abuse. \emph{Journal of Abnormal
Psychology, 114}(2), 259--268.
\url{https://doi.org/10.1037/0021-843X.114.2.259}

Ullman, S. E., \& Peter-Hagene, L. C. (2014). Social reactions to sexual
assault disclosure, coping, perceived control and PTSD symptoms in
sexual assault victims. \emph{Journal of Community Psychology, 42}(4),
495--508. \url{https://doi.org/10.1002/jcop.21624}

Vigdal, J. S., \& Brønnick, K. (2022). A systematic review of
``helicopter parenting'' and its relationship with anxiety and
depression. \emph{Frontiers in Psychology, 13}, Article 872981.
\url{https://doi.org/10.3389/fpsyg.2022.872981}
\end{apareferences}

\end{document}
